	
	% ==========================
	% Abstract
	% ==========================
	\clearpage
	\hypersetup{pageanchor=true}
	\pagestyle{summary}
	
	\vspace*{\fill}
	\begin{center}{\Large \textbf{\ThesisHeadingText{abstract}}}\end{center}
	\vspace{0.5cm}
	This work presents a theoretical and numerical study of exact solutions of the Einstein field equations, with particular emphasis on the Schwarzschild and Kerr metrics. After recalling the geometric foundations of general relativity, the geodesic structure of these spacetimes is analyzed, including the innermost stable circular orbit and its dependence on the black hole spin. A numerical ray-tracing procedure is then developed to integrate null geodesics and compute gravitational lensing effects and the black hole shadow, whose predicted size is compared with Event Horizon Telescope observations of M87*. Finally, the dynamical response of these spacetimes to small perturbations is investigated through the quasinormal mode spectrum, providing a black hole spectroscopy framework that is illustrated using the ringdown of the gravitational-wave event GW150914. The consistency between the analytical, numerical, and observational results obtained throughout this work supports the Kerr metric as the standard description of astrophysical black holes.
	\vspace{0.3cm}\\
	\noindent\textbf{\ThesisHeadingText{keywords}} General relativity, Schwarzschild metric, Kerr metric, black hole, geodesics, gravitational lensing, quasinormal modes.
	\vspace*{\fill}
	
	\newpage
	
